\section{The FK sphere and its quantum area}
\label{finite:section}

Throughout this Section, \(q\in(0,4)\) is fixed, with
\(\gamma\in(\sqrt2,2)\) determined by Equation~\eqref{eq:q-gamma}. Write
\[
 \delta_n=\frac{d_n}{B(2n)},
\]
where \(B\) is the deterministic distance denominator constructed in
Lemma~\ref{comp:polynomial-algebra}.  In particular,
\(B(2n)\to\infty\).  Distances in this Section always use all the primal
edges of \(M_n\). Equip its vertices with
\[
                      \mu_n(v)=\frac{\deg(v)}{2n},
\]
where a loop contributes twice to the degree. Our goal is convergence
of this metric probability space to the ordinary unit-area sphere.

We first specify the inputs from the preceding Sections.  This also
separates the local comparison theorem from the compactness argument
which remains to be proved here.
\begin{enumerate}[label=(F\arabic*)]
\item
For every sequence of time units tending to infinity, there is a
subsequence along which the bilateral passage rule has the comparisons
of Corollary~\ref{comp:bilateral-comparison}, with a
\emph{deterministic} constant
\[
                 c\in[c_-,c_+]\subset(0,\infty).
\]
The same \(c\) applies to all the intrinsic local charts in that
extraction.  These comparisons include lower bounds for every retained
path and upper bounds for following a prescribed rectifiable continuum
path with endpoint slack.  On subpaths compactly contained in a
buffered chart, the comparison is with the intrinsic LQG metric of that
chart, so that it is determined by the real field there.
\item
The endpoint joining conclusion of
Corollary~\ref{spike:endpoint-modulus} holds in these units.  It applies to
arbitrary discrete endpoints whose projected locations approach one
another in a compact chart, with any prescribed open clearance.
\item
The finite-conditioned transfer
Theorem~\ref{loc:finite-transfer}, in its spatial form
Theorem~\ref{loc:spatial-kernels}, preserves the conditional local
kernels on strict interior charts, with the same time and distance
units.  Its observables include the universal endpoint tests in (F2).
We use the actual-field absolute continuity of
Lemma~\ref{bg:sphere-field-ac} to transfer their metric support.
\end{enumerate}
None of these inputs asserts tightness of the whole finite metric
space.  In particular, (F1) concerns passages with endpoint slack,
whereas (F2) is needed to start paths at specified lattice vertices.

\subsection{The conditioned surface and its local metric}

Lemma~\ref{bg:finite-law} identifies the inventory word conditioned
to reduce to the empty word with the required FK-weighted map. We now
identify the surface encoded by its limit and transfer the local metric
statements to that surface.

By \cite[Theorem 1.8]{GwynneSunFinite}, the resolved contour pair, with
time divided by \(2n\) and heights by \(\sqrt{2n}\), in the fixed
boundary-length convention of Section~\ref{sec:framework}, converges to the duration-one correlated
quadrant excursion.  Its correlation is
\[
 \rho=-\cos(\pi\gamma^2/4).
\]
By \cite[Theorem 1.1]{MillerSheffieldSphere}, this excursion determines
the ordinary unit-area DMS sphere and its space-filling curve
\[
                 (\mathcal S,h,\eta),\qquad \eta:[0,1]\to\mathcal S.
\]
Here \(\eta(0)=\eta(1)=p_*\), and
\[
           \mu_h(\eta([s,t]))=t-s,\qquad 0\le s\le t\le1.
\]
We forget all distinguished points except the point \(p_*\) used
temporarily in the proof.  The quantum-surface law is the ordinary unit-area law.
Additional area marks below serve only to choose coordinate
representations; they do not change this law or reroot the exploration.

Consider an arbitrary subsequence of positive integers \(n\).
Use (F1) to select a further subsequence and its deterministic
constant \(c\), and extract the finite chart data jointly along it.
We work with a coupling in which the contours converge uniformly and
the countable chart tests converge.  Whenever a joint coupling is
refined below, we retain these tests and a countable family of their
local modulus and comparison error variables.  Errors converging to
zero in probability can first be made almost surely convergent by a
diagonal subsequence.  Thus subsequent joint extractions preserve the
local conclusions.  All subsequent subsequences in this Section are
taken from this selected sequence.

For each vertex \(v\), choose one of its contour representatives
\(t_n(v)\in[0,1]\), and set
\[
                    \phi_n(v)=\eta(t_n(v)).
\]
For the root vertex choose its representative at time \(0\).
The choices of representatives are immaterial asymptotically:
all the representatives of a primal vertex have projected diameter
tending uniformly to zero.  The same is true of the two endpoints of
a primal edge.  Here the number of encoding relations needed is
bounded independently of the edge and of \(n\).
Two representatives of a primal vertex satisfy the first-coordinate
contour-tree identity.  A tree edge is represented by adjacent tour
times.  For a remaining primal edge, the two triangle steps of its
quadrangle are matched in the second coordinate: use their horizontal
matching relation and, if necessary, one adjacent-time move at each
end to reach the primal endpoint visits.  Finally one may change to
the chosen representatives of those two primal vertices by two
first-coordinate identities.  This uses at most five relations.

Consequently a violating sequence would give, after extraction of
these boundedly many times, a limiting horizontal chain in which each
adjacent-time move has disappeared.  The horizontal-identification
rule identifies the endpoints, so its image under \(\eta\) is a
single point, a contradiction.  This argument uses the resolved
tree/cotree sewing, not a graph path of an uncontrolled number of
edges.  Thus, for any fixed auxiliary metric
\(\sigma\) inducing the topology of \(\mathcal S\),
\begin{equation}
 \max_{\{v,w\}\in E(M_n)}
       \sigma(\phi_n(v),\phi_n(w))\longrightarrow0 .
 \label{finite:edge-mesh}
\end{equation}
The original root vertex projects to \(p_*\).
We can interpolate the projected edges by arcs of vanishing
\(\sigma\)-diameter.  Confinement below allows these vanishing
interpolation errors.
For open \(U\subset\mathcal S\), the notation
\(\delta_n(v,w;U)\) denotes the infimum of normalized lengths of
graph paths from \(v\) to \(w\) whose vertex projections lie in \(U\).
All uses have open clearance, so the edge interpolation convention
does not change the limiting statements.

\begin{lemma}[Spatial transfer to the ordinary sphere]
\label{finite:spatial-field-transfer}
In the preceding joint extraction, the bilateral intrinsic local
comparisons and endpoint tests in (F1)--(F2) hold in the finite sphere
on compact patches away from $p_*$, with the same deterministic
constant $c$.
\end{lemma}
\begin{proof}
Adjoin two auxiliary points $Q_1,Q_2$, sampled independently from
$\mu_h$ conditional on the surface and independently of the exploration
and chart observations given that surface and its contour.  They may
be realized as $\eta(T_1),\eta(T_2)$ for independent uniform times
adjoined to the joint coupling.  By the area-marking property
\cite[Proposition~A.13]{DMS} and the three-mark representation in
\cite[Definition~2.2 and Lemma~2.3]{BorgaGwynneSunPermutons},
$(\mathcal S,Q_1,Q_2,p_*)$ has the ordinary three-marked unit-area
sphere law.  Represent it in $\mathbb C$ with the marks at $0,1,\infty$.
The identification with the DMS sphere includes the field, as explained
in Lemma~\ref{bg:sphere-field-ac}.

Before applying the quantum-area clock, the space-filling curve in
\cite[Theorem~1.1]{MillerSheffieldSphere} has the whole-plane
space-filling SLE law from infinity to infinity, independent of the
real field.  Construct it from an imaginary field modulo its period
with stationary phase.  This is the same directed imaginary-field and
curve law as in the cone experiment.  Conformal invariance, including
the imaginary-field phase transformation, preserves independence under
the affine coordinate normalizations used here.

Fix $K\Subset U$, where $U$ is a relatively compact disk in
$\mathbb C\setminus\{0,1\}$, with the buffers needed by
Theorem~\ref{loc:spatial-kernels}.  Let $\mu_{\rm sph}$ and
$\mu_{\rm cone}$ be the laws of the actual real-field restrictions
on $U$, using the cone representation in
Lemma~\ref{bg:sphere-field-ac}.  Let $\nu$ be the law of the entire
directed imaginary-field and space-filling-curve object, and let
$\Pi_h$ be the common conditional law of the area marks in $U$.
The field absolute continuity and the common mark kernel give
\begin{equation}\label{finite:field-product-ac}
 \mu_{\rm sph}(dh)\,\nu(d\vartheta)\,\Pi_h(d\pi)
 \ll
 \mu_{\rm cone}(dh)\,\nu(d\vartheta)\,\Pi_h(d\pi).
\end{equation}
The entire directed imaginary-field object retains the global order
of its local traversals.  The finite order needed by a selected chart,
its locally clocked inputs, and the relevant mark data are measurable
functions of the displayed variables.  Thus the order variable in
Theorem~\ref{loc:spatial-kernels} is retained in this comparison.

That theorem gives the same conditional local observation kernel
in the two experiments.  The intrinsic metric on $U$ is determined
by $h|_U$.  Hence failure of a retained intrinsic local comparison or
endpoint test is a measurable event of the local input and observation.
On the cone its conditional probability is zero almost everywhere
for the right-hand input law in
Equation~\eqref{finite:field-product-ac}.  Absolute continuity makes
it zero for the sphere input law, and the common-kernel identity proves
the test there.  The transported metric is the internal metric on $U$;
an unrestricted metric depending on the exterior is not used in this
step.  Intersect over a countable disk cover, the protected-chart
atlas, endpoint tests and rational open losses.  We obtain the local
statements away from $p_*,Q_1,Q_2$.  Their finite kernels apply because
all contour preimages of a compact set avoiding $p_*$ lie in a strict
interior time interval.

Adjoin a second independent pair $Q_3,Q_4$ of area points, retaining
the same sphere, exploration, discrete data and reference extraction.
Apply the argument to each three-marked marginal separately, in the
representations $(Q_1,Q_2,p_*)=(0,1,\infty)$ and
$(Q_3,Q_4,p_*)=(0,1,\infty)$; no field absolute continuity conditional
on the other pair is asserted.  The four auxiliary points are distinct
and avoid $p_*$ almost surely, since $\mu_h$ is diffuse.  The two
coordinate changes differ by a complex affine map.  Quantum-area
clocks, quantum lengths and the chart observations have the intrinsic
covariance of Theorem~\ref{loc:spatial-kernels}; the internal LQG
metrics have conformal covariance.  Therefore the two conclusions
are statements about the same intrinsic comparisons, with the same
chosen $c$.  The exploration is not rerooted, the quantum units are
unchanged, and no further bilateral normalization is selected.

The union of the two off-mark covers is
$\mathcal S\setminus\{p_*\}$, because the auxiliary mark sets are
disjoint.  Compact confinements admit finite subcovers, proving the
lemma.
\end{proof}

\begin{lemma}[Local finite-volume metric statements]
\label{finite:local-metric}
In the coupling above, the following statements hold on
\(\mathcal S\setminus\{p_*\}\).
\begin{enumerate}[label=(\roman*)]
\item
Let \(P\) be a \(D_h\)-rectifiable path whose image is a compact subset
of an open set \(U\subset\mathcal S\setminus\{p_*\}\).
If \(\phi_n(v_n)\to P(0)\) and \(\phi_n(w_n)\to P(1)\), then
\[
 \limsup_{n\to\infty}
     \delta_n(v_n,w_n;U)
       \le c\,\operatorname{len}_{D_h}(P).
\]
An arbitrarily small open enlargement of the image of \(P\) may be
used in place of \(U\).
\item
For paths confined to a compact subset of one buffered chart, every
limiting endpoint pair \(x,y\) satisfies
\[
       \liminf_{n\to\infty}
          \operatorname{len}_{\delta_n}(P_n)
                  \ge cD_h(x,y).
\]
This assertion holds for all choices of the paths and endpoints.
\item
The local endpoint joining modulus is uniform over all discrete
endpoints in a compact chart, with paths confined to any specified
open enlargement.
\end{enumerate}
\end{lemma}

\begin{proof}
A compact confinement away from \(p_*\) has all its contour preimages
in a strict interior time interval: otherwise a sequence of preimages
approaching \(0\) or \(1\), together with continuity of \(\eta\), would
put \(p_*\) in that confinement.  We may therefore use the finite
transfer (F3) on a countable collection of buffered charts covering
the confinement.

Lemma~\ref{finite:spatial-field-transfer} transfers the comparisons
in (F1) and the endpoint tests in (F2), using the common spatial
kernels and absolute continuity of the actual local fields.  The
transferred constant is the same deterministic \(c\) on every chart.
An intrinsic lower bound implies the corresponding ambient lower
bound, since an internal metric dominates the ambient metric.  The
length of a fixed path compactly contained in a chart is unchanged
when computed using its internal metric.

The lower comparison and the universal endpoint tests give
(ii) and (iii) directly.  For (i), first divide \(P\) into finitely
many subpaths lying with clearance in buffered charts.  Follow these
subpaths using (F1).  Their endpoints initially have only spatial
slack.  Part (iii) joins these endpoints and joins the first and last
ones to \(v_n,w_n\), at a total cost tending to zero.  One may keep
every such join inside the prescribed enlargement.  Sum the
subpath lengths, and then remove the finite partition and cost
slack.  A countable atlas and the universal endpoint observables
make all three conclusions simultaneous.  In particular, the last
step does not infer a universal endpoint assertion from existence
of a limiting path certificate.
\end{proof}

\begin{lemma}[Lower bounds on compact punctured routes]
\label{finite:compact-lower}
Let \(K\subset\mathcal S\setminus\{p_*\}\) be compact.
Suppose that \(P_n\) has projected image in \(K\), up to vanishing
errors, and that its endpoint projections converge to \(x,y\).
Then
\[
       \liminf_{n\to\infty}
          \operatorname{len}_{\delta_n}(P_n)
                    \ge cD_h(x,y).
\]
\end{lemma}

\begin{proof}
It suffices to consider a subsequence on which the left side is
finite and the path costs are bounded by some \(L<\infty\).
Take a compact neighborhood \(K'\) of \(K\) avoiding \(p_*\).
Choose finitely many buffered charts with compact cores whose
interiors cover \(K'\).
For sufficiently small \(\rho>0\), every
\(\sigma\)-ball of radius \(3\rho\) centered in \(K'\) lies in one
of these cores.

There is \(a_\rho>0\) such that, for all sufficiently large \(n\),
a retained path in the fixed compact core of one of these buffered
charts, with endpoints in \(K'\) making a \(\sigma\)-advance of at
least \(\rho/2\), has cost at least \(a_\rho\).
Otherwise one could select violating paths with costs tending
below half of
\[
 c\min\{D_h(u,v):u,v\in K',\ \sigma(u,v)\ge\rho/2\}>0,
\]
extract their endpoints, and contradict
Lemma~\ref{finite:local-metric}(ii).
The minimum is positive because \(D_h\) is a metric inducing the
compact topology on \(K'\).

Starting at the initial vertex of \(P_n\), stop at the first vertex
whose projection is at least \(\rho\) away from the preceding stop, and
repeat.  Equation~\eqref{finite:edge-mesh} absorbs the overshoots.
Every complete segment lies in a buffered chart and costs at least
\(a_\rho\).  There are therefore at most \(L/a_\rho\) complete
segments.  After a further extraction their endpoints converge.
Apply the local lower bound to the finitely many complete
segments and sum.  The terminal incomplete segment has endpoint
\(D_h\)-distance at most
\[
 \omega_{K'}(2\rho)
   :=\sup\{D_h(u,v):u,v\in K',\ \sigma(u,v)\le2\rho\}.
\]
The triangle inequality gives
\[
 cD_h(x,y)
 \le \liminf_n\operatorname{len}_{\delta_n}(P_n)
       +c\omega_{K'}(2\rho).
\]
Uniform continuity gives \(\omega_{K'}(2\rho)\to0\), proving the
claim.
\end{proof}

\subsection{Small surrounding bands}

We record the continuum fact used to bypass the point \(p_*\).
An internal diameter is always taken in the indicated open
neighborhood of the set.

\begin{lemma}[Vanishing surrounding bands]
\label{finite:small-bands}
Almost surely there are compact annular bands \(K_m\) and open
neighborhoods \(U_m\), with
\[
 K_m\subset U_m\subset\mathcal S\setminus\{p_*\},
 \qquad
 \sup_{z\in U_m}D_h(z,p_*)\longrightarrow0,
\]
such that each \(K_m\) contains, with positive topological clearance,
a separator between \(p_*\) and the complement of a neighborhood
shrinking to \(p_*\), and
\[
             \operatorname{diam}_{D_h(\cdot,\cdot;U_m)}K_m
                          \longrightarrow0.
\]
The bands can be chosen with an inner and an outer collar, so that
paths crossing them are stable under vanishing perturbations.
\end{lemma}

\begin{proof}
First work in a conformal chart on a compact set avoiding the field
marks.  The internal small-box connection estimate in
Theorem~\ref{bg:metric-inputs} gives a
random upper bound tending to zero for the internal diameter of a
small box in a fixed bounded enlargement.  The bound is uniform
over boxes in this compact set.  Cover a square annulus by a fixed
number of overlapping boxes, using enlargements which avoid its
central point.  Chaining through the overlaps gives an internal
diameter tending to zero for the annulus.  Thinner concentric
square annuli provide both collars.  Since the box estimate holds
simultaneously on compacts, this construction applies to every
point of the chart.  A countable chart cover gives the property at
every point except the finitely many field marks.

The property just established says that the point has shrinking
surrounding bands with vanishing internal joining diameter.  It is
a property of the metric surface at the point, and is independent
of which auxiliary marks were used to embed the field.  The
quantum area measure is diffuse, so the property holds at
\(\mu_h\)-almost every point of the unmarked surface.
The area-rerooting property of the ordinary sphere says that the
distinguished point \(p_*\), conditionally on that surface, is
sampled from its area measure; see
\cite[Remark 1.3]{MillerSheffieldSphere}, citing
\cite[Proposition A.13]{DMS}.
It follows that the property holds at \(p_*\).
Finally, the metric induces the surface topology, so the bands
and their sufficiently small collars have
\(\sup_{z\in U_m}D_h(z,p_*)\to0\).
By taking strict bounds and collars, all choices can be made
measurably from countable chart and rational-box configurations.
\end{proof}

\begin{remark}
\label{finite:band-input}
The internal connection estimate in this proof is essential.
The conclusion of Lemma~\ref{finite:small-bands} is not being
deduced solely from geodesicity and the topology of the sphere.
The rerooting argument transfers a property already proved on
ordinary local field charts; it does not require a uniform
Radon--Nikodym bound at a shrinking neighborhood of a field mark.
\end{remark}

\begin{proposition}[Uniform punctured distances]
\label{finite:punctured}
For every compact \(K\subset\mathcal S\setminus\{p_*\}\),
\[
 \sup_{\substack{v,w\in V(M_n)\\
                 \phi_n(v),\phi_n(w)\in K}}
 \left|\delta_n(v,w)-cD_h(\phi_n(v),\phi_n(w))\right|
                       \longrightarrow0 .
\]
The distances on the left are unrestricted graph distances.
An empty supremum of nonnegative errors is interpreted as zero.
\end{proposition}

\begin{proof}
Fix sequences of endpoints with projections converging to
\(x,y\in K\).  For the upper bound take a \(D_h\)-geodesic from
\(x\) to \(y\).
If it meets a sufficiently small hole around \(p_*\), replace its
portion between first entrance and last exit of a surrounding
band by a path in \(U_m\).  The additional length tends to zero
by Lemma~\ref{finite:small-bands}.  The resulting path avoids
\(p_*\), is compactly contained in its complement, and has length
at most \(D_h(x,y)+o_m(1)\).
Lemma~\ref{finite:local-metric}(i) gives
\[
            \limsup_n\delta_n(v_n,w_n)\le cD_h(x,y).
\]

For the lower bound, choose paths from \(v_n\) to \(w_n\) with
costs within \(1/n\) of \(\delta_n(v_n,w_n)\).
The upper bound makes their costs bounded.
Fix a small surrounding band with collars.  If a path enters
the inner hole, retain only its part up to its first entrance
and its part after its last exit.  The retained parts lie in a
compact punctured region, and the entrance and exit vertices
lie in a small neighborhood of \(p_*\), with an error vanishing
by Equation~\eqref{finite:edge-mesh}.
Apply Lemma~\ref{finite:compact-lower} separately to these two
parts.  The omitted cost is nonnegative.  The \(D_h\)-distance
between the two limiting entrance points tends to zero as the
band shrinks, by the triangle inequality through \(p_*\).
If the path never enters the inner hole, apply
Lemma~\ref{finite:compact-lower} to the whole path.
In both cases,
\[
           \liminf_n\delta_n(v_n,w_n)\ge cD_h(x,y).
\]
The collars ensure that no edge can skip the band.

This proves convergence for arbitrary sequences of endpoints
with limiting locations in \(K\).  If the asserted uniform
convergence failed, compactness of \(K\) would supply violating
endpoint sequences with convergent projections.  The result
just proved would contradict that violation.
\end{proof}

\subsection{Sample arrays and the two roots}

Punctured convergence leaves uncontrolled only those vertices whose
projections approach the exploration root. We use two conditionally
independent uniform corner roots of the same map, together with one
common sequence of corner samples. The sample arrays identify the two
continuum completions, and the roots become distinct points in that
completion. A vertex escaping every finite sample net would have to
approach both roots. Separating bands determine its distance differences
to the samples and force the two root points to coincide.

Conditional on the unrooted decorated map, let \(R_n^0\) be its
original root corner, let \(R_n^1\) be a second independent
uniform corner, and let
\[
                     X_n^1,X_n^2,\ldots
\]
be further independent uniform corners.  Each corner is
represented in the metric by its incident primal vertex.
Reroot at \(R_n^1\) and form the second inventory exploration.

A corner corresponds to its outgoing dart, so an \(n\)-edge map has
\(2n\) corners. Exactly \(\deg(v)\) are incident to a vertex \(v\),
including two for each loop at \(v\). Thus each sampled corner,
represented by its vertex, has the probability law \(\mu_n\) defined
above. This remains true for a one-vertex map with loops.
Conditional on the unrooted decorated map,
the original root and these additional corners are
exchangeable.  Indeed the map weight does not depend on its
root.  Choosing a uniform corner in a fixed representative
gives the uniform rooted isomorphism class: the automorphism
group acts freely on oriented corners, so all its corner
orbits have the same size.  The subsequent independent
choices preserve exchangeability of every metric distance
array.

Both explorations have the same marginal law.  Extract their
contours and local tests jointly, using the same normalization
\(B(2n)\) and the same bilateral reference extraction.  They
give
\[
 (\mathcal S^j,cD^j,\mu^j,\eta^j,p_j),\qquad j=0,1.
\]
The discrete samples retain their common labels in the two
explorations.

\begin{lemma}[Common dense sample array]
\label{finite:sample-array}
After a joint extraction, for each \(j\) the locations
\(X_n^i\) converge in exploration \(j\) to points \(X^{j,i}\)
which are iid from \(\mu^j\), conditionally on that surface and
its curve.  For every \(k<\infty\),
\[
 \left(\delta_n(X_n^i,X_n^\ell)\right)_{i,\ell\le k}
 \longrightarrow
 \left(cD^j(X^{j,i},X^{j,\ell})\right)_{i,\ell\le k}.
\]
The labeled arrays identify the two limiting compact metric
probability spaces by a unique measure-preserving isometry sending
\(X^{0,i}\) to \(X^{1,i}\) for every \(i\).
\end{lemma}

\begin{proof}
The triangle tour passes through every Tutte edge once, and
Tutte edges correspond to corners of the primal map.
Consequently an independent uniform corner is a uniform tour
time, with only the fixed endpoint/triangle convention of the
encoding.  Its normalized index is uniform on the discrete
time grid.  Conditional on an exploration, the indices of
independently sampled corners are independent uniform
indices.  They converge jointly with the contours to iid
uniform times \(U_i^j\), independent of the limiting contour
data.  Thus
\[
                         X^{j,i}=\eta^j(U_i^j).
\]
Area parametrization gives the asserted conditional area law.

Every fixed sample avoids \(p_j\) almost surely.  Apply
Proposition~\ref{finite:punctured} to the finitely many
locations in a fixed array, and then use a diagonal
extraction over \(k\).  The numerical arrays obtained from
the two explorations are identical, since their discrete
entries were identical.

Each area measure has full support. The marginal sample sequences are
therefore dense almost surely, and Lemma~\ref{measure:common-completion}
extends the common labeled array to a unique measure-preserving isometry
of the two compact spaces. This uses the two marginal sampling laws;
no conditional iid assertion given both explorations, or prior
convergence of metric-measure spaces, is needed.
\end{proof}

Use this isometry to regard the limits as one compact metric
probability space \((S,d,\mu_h)\), where \(d=cD_h\), and write
\(X^i\) for the common labeled samples.  We next include the root labels in
the limiting arrays.  Their finite arrays are tight by
exchangeability with sample arrays, so they can be included
in the same diagonal extraction.

\begin{lemma}[Root completion points]
\label{finite:root-completion}
Each root label determines a point \(Z_j\in S\) such that
\[
              \delta_n(R_n^j,X_n^i)\longrightarrow d(Z_j,X^i)
              \quad\text{for every }i.
\]
Their mutual distance also converges to \(d(Z_0,Z_1)\).
Moreover \(Z_0\ne Z_1\) almost surely.
\end{lemma}

\begin{proof}
Let \(l_{j,i}\) denote the limiting distance from root \(j\)
to sample \(i\).  Exchangeability gives the same law to every
finite array including the root labels as to a sample array
with the corresponding number of additional labels.
It follows that
\begin{equation}
                       \inf_{i\ge1}l_{j,i}=0
                       \quad\text{almost surely}.
 \label{finite:root-near-sample}
\end{equation}
One can see this without reconstructing a limiting measure:
for fixed \(k\), the law of
\(\min_{i\le k}l_{j,i}\) is the law of the distance from one
iid area sample to the first \(k\) other samples.  These
distances decrease to zero by their almost sure density.

Choose indices \(i_m\) with \(l_{j,i_m}\to0\).  The triangle
inequalities in the finite arrays pass to the limit, giving
\[
 d(X^{i_m},X^{i_\ell})\le l_{j,i_m}+l_{j,i_\ell}.
\]
Thus \(X^{i_m}\) is Cauchy and has a limit \(Z_j\in S\).
The same inequalities imply
\[
 |l_{j,i}-d(X^{i_m},X^i)|\le l_{j,i_m},
\]
so \(l_{j,i}=d(Z_j,X^i)\) for all \(i\).
This also proves uniqueness.  If \(l_{01}\) is the limiting
root-to-root distance, then
\[
                       |l_{01}-l_{1,i}|\le l_{0,i}.
\]
Taking samples approaching \(Z_0\) gives
\(l_{01}=d(Z_0,Z_1)\).

Finally, by exchangeability \(l_{01}\) has the law of the
distance between two independent area samples in
\((S,d,\mu_h)\).  The area measure is diffuse, so this distance
is strictly positive almost surely.
\end{proof}

We need a simple deterministic observation about a separating
set in a graph.
\begin{lemma}[Separation and distance profiles]
\label{finite:profile}
Let \((V,\delta)\) be a finite graph with its length metric.
Let \(C\subset V\) meet every path from \(v\) to each of the
vertices \(x_1,\ldots,x_k\), and put
\(e=\delta(v,C)\).  Then
\[
 e+\inf_{z\in C}\delta(z,x_i)
 \le \delta(v,x_i)
 \le e+\sup_{z\in C}\delta(z,x_i).
\]
In particular, if
\(\sup_{z\in C}|\delta(z,x_i)-a_i|\le t\) for each \(i\), then
\[
 |(\delta(v,x_i)-\delta(v,x_1))-(a_i-a_1)|\le2t .
\]
The last bound is independent of \(e\).
\end{lemma}

\begin{proof}
A shortest path from \(v\) to \(x_i\) first meets \(C\) at
some vertex \(z\).  Its two portions have costs at least
\(e\) and \(\delta(z,x_i)\), giving the lower bound.
For the upper bound, join \(v\) to a point minimizing its
distance to \(C\), then continue to \(x_i\).  Subtract the
two resulting intervals to obtain the last assertion.
\end{proof}

For each exploration \(j\), the bands of
Lemma~\ref{finite:small-bands}, their collars, and
Equation~\eqref{finite:edge-mesh} produce discrete vertex sets
\(C_{n,m}^j\) with the following properties.
For fixed \(m\), vertices projecting sufficiently close to
\(p_j\) are separated by \(C_{n,m}^j\) from every fixed sample
outside the band.  For all large \(m\), every fixed finite
list of samples is outside.  Moreover, by
Proposition~\ref{finite:punctured},
\begin{equation}
 \lim_{m\to\infty}\limsup_{n\to\infty}
 \sup_{z\in C_{n,m}^j}
       |\delta_n(z,X_n^i)-d(p_j,X^i)|=0
 \quad\text{for every fixed }i .
 \label{finite:band-profile}
\end{equation}
Here and below \(p_j\) is viewed in \(S\) by the
label-preserving isometry.
To construct these sets, take the vertices whose projections
lie in a closed thick band around a separator compactly
inside \(K_m\).  A projected path from its inner side to its
outer side has a vertex in that thick band once the edge
mesh is smaller than its collars.  The band remains a
compact punctured set for fixed \(m\), so
Proposition~\ref{finite:punctured} applies uniformly there.
Its continuum distance profile tends to that of \(p_j\)
because its points converge uniformly to \(p_j\).

\begin{lemma}[Identification of the root with the missing point]
\label{finite:root-identification}
In the common completion, \(Z_j=p_j\), for \(j=0,1\).
Consequently \(p_0\ne p_1\) almost surely.
\end{lemma}

\begin{proof}
Set
\(e_{n,m}^j=\delta_n(R_n^j,C_{n,m}^j)\).
For any fixed sample \(i\), all sufficiently small bands
separate the root from that sample.  Therefore
\[
               e_{n,m}^j\le\delta_n(R_n^j,X_n^i).
\]
First take a limsup in \(n\), then in \(m\), and finally the
infimum over \(i\).  Equation~\eqref{finite:root-near-sample}
gives
\[
                 \lim_{m\to\infty}\limsup_n e_{n,m}^j=0.
\]
Apply Lemma~\ref{finite:profile} with \(v=R_n^j\), and use
Equation~\eqref{finite:band-profile}.  For every \(i\),
\[
                   d(Z_j,X^i)=d(p_j,X^i).
\]
The sample set is dense, so \(Z_j=p_j\).
The last assertion follows from
Lemma~\ref{finite:root-completion}.
\end{proof}

\begin{proposition}[Tight finite metric nets]
\label{finite:metric-nets}
For every \(\varepsilon>0\),
\[
 \lim_{k\to\infty}\limsup_{n\to\infty}
 \mathbb P\!\left[
   \max_{v\in V(M_n)}\min_{1\le i\le k}
       \delta_n(v,X_n^i)>\varepsilon
 \right]=0
\]
along the selected subsequence.
\end{proposition}

\begin{proof}
Suppose otherwise.  There are \(\varepsilon,\theta>0\),
integers \(k_n\to\infty\), and a further subsequence on which
the displayed event with \(k=k_n\) has probability at least
\(\theta\).  On this event choose a violating vertex \(v_n\)
by a fixed finite ordering; elsewhere choose any vertex.
Include the event indicator and contour representatives of
\(v_n\) in both explorations in the joint extraction.
These additional variables are tight: the indicator lies
in \(\{0,1\}\) and the representative times lie in \([0,1]\).
We do not condition the laws of either limiting sphere on
the event.  In the resulting coupling the indicators
converge almost surely to a variable \(I\in\{0,1\}\).
Put \(E=\{I=1\}\).  Then \(\mathbb P(E)\ge\theta\), and on
\(E\) the selected vertices violate the net condition for
all sufficiently large \(n\).

On \(E\), the projected vertex must converge to \(p_j\) in
each exploration.  Indeed, a limiting location different
from \(p_j\) lies in a compact punctured region.  A fixed
sample can be chosen arbitrarily close to that location.
Proposition~\ref{finite:punctured} would then give
\(\delta_n(v_n,X_n^i)<\varepsilon\) for some fixed \(i\) and
all sufficiently large \(n\), contrary to \(k_n\to\infty\)
and the choice of \(v_n\).

Fix \(i\).  In either exploration, for each sufficiently
large fixed \(m\), all paths from \(v_n\) to \(X_n^i\) and
\(X_n^1\) meet \(C_{n,m}^j\), eventually on \(E\).
Equations~\eqref{finite:band-profile} and
Lemma~\ref{finite:profile} imply
\[
 \delta_n(v_n,X_n^i)-\delta_n(v_n,X_n^1)
       \longrightarrow d(p_j,X^i)-d(p_j,X^1).
\]
The possibly unbounded escape cost from \(v_n\) cancels.
The numerical expression on the left is identical in both
explorations, so on \(E\)
\[
 d(p_0,X^i)-d(p_0,X^1)
       =d(p_1,X^i)-d(p_1,X^1)
       \qquad\text{for every }i.
\]
Density and continuity imply that
\(d(p_0,x)-d(p_1,x)\) is constant for \(x\in S\).
At \(x=p_0\) its value is \(-d(p_0,p_1)\), whereas at
\(x=p_1\) its value is \(d(p_0,p_1)\).
Thus \(p_0=p_1\) on \(E\), contradicting
Lemma~\ref{finite:root-identification}.
\end{proof}

\begin{corollary}[The subsequential metric and area limit]
\label{finite:subsequential}
Along the selected subsequence,
\[
                (V(M_n),\delta_n,\mu_n)
                      \ \Longrightarrow\ (S,cD_h,\mu_h)
\]
in Gromov--Hausdorff--Prokhorov topology.
Every original subsequence has a further subsequence of this
form, with \(c\) deterministic and in \([c_-,c_+]\).
\end{corollary}

\begin{proof}
The samples \(X_n^i\), represented by their incident vertices, are
conditionally iid with law \(\mu_n\).
Lemma~\ref{finite:sample-array} gives convergence of every fixed
finite distance array to the array of conditionally iid area samples
in \((S,cD_h,\mu_h)\). The limiting area measure has full support.
Proposition~\ref{finite:metric-nets} gives precisely the covering
condition in Lemma~\ref{measure:criterion}. Applying that lemma proves
the GHP assertion with the identified quantum-area measure.

The constant \(c\) was selected in (F1) before either finite exploration
was extracted. It is deterministic and common to all charts. Neither
sampling, completion nor the measured convergence criterion changes
this multiplier. The construction began with an arbitrary subsequence,
which proves the last assertion.
\end{proof}

\subsection{A deterministic normalization for all integers}

We apply the common normalization lemma to the measured limit just
proved. Let
\[
 H_n=\operatorname{diam}(V(M_n),d_n),
 \qquad H=\operatorname{diam}(S,D_h),
\]
and let \(m\) be the center of Lemma~\ref{measure:center}.
The variable \(H\) is finite and strictly positive almost surely.
Define the deterministic numbers
\begin{equation}\label{finite:normalization}
 b_n=\exp\!\left(
       m\big(\log(H_n\vee1)\big)-m(\log H)
      \right),
 \qquad a_n(q)=b_n^{-1}.
\end{equation}
The truncation includes maps with a single vertex and loops, whose
unscaled graph diameter is zero.

\begin{proposition}[The full-sequence FK limit]
\label{finite:full-sequence}
The numbers in Equation~\eqref{finite:normalization} satisfy
\mbox{\(a_n(q)>0\)}, \(a_n(q)\to0\), and
\[
              (V(M_n),a_n(q)d_n,\mu_n)
                    \ \Longrightarrow\ (S,D_h,\mu_h)
\]
in Gromov--Hausdorff--Prokhorov topology as \(n\to\infty\) through
all positive integers.
\end{proposition}

\begin{proof}
Set \(s_n=B(2n)\), which tends to infinity by
Lemma~\ref{comp:polynomial-algebra}.
Corollary~\ref{finite:subsequential} supplies every hypothesis of
Lemma~\ref{measure:normalization}, including the deterministic nature
of the subsequential multiplier. The normalization prescribed by that
lemma is exactly Equation~\eqref{finite:normalization}. Its conclusion
proves the proposition, with the quantum-area measure unchanged.
This proves the critical FK clause of Theorem~\ref{thm:main}.
\end{proof}

The normalization uses unconditional diameter laws and is deterministic.
The proof uses the deterministic multiplier assertion in (F1); it does
not replace a random subsequential multiplier by a deterministic one.
All constants and all subsequence choices in this Section may depend
on the fixed value of \(q\). No endpoint statement or uniformity as
\(q\) approaches \(0\) or \(4\) is used.
